\section{Introduction}

Affective computing systems increasingly estimate a person's emotional state
from facial appearance. Two failure modes make naive deployment irresponsible:
(i) models are overconfident, reporting a precise emotion even when the input is
ambiguous, occluded, or out of distribution; and (ii) the dimensional structure
of affect is flattened into a handful of discrete categories that discard
intensity and ambiguity. We address both by predicting \emph{continuous} valence
and arousal (VA) on the circumplex model of affect~\cite{russell1980circumplex}
and, crucially, by attaching a \emph{calibrated} measure of uncertainty to every
prediction, with an explicit abstention path when the signal cannot support a
confident estimate.

\paragraph{Contribution.} Our contribution is calibration and abstention, not raw
accuracy. Concretely: (1) a deep evidential regression head~\cite{amini2020evidential}
that produces aleatoric and epistemic uncertainty for VA in one forward pass;
(2) a split conformal layer~\cite{vovk2005algorithmic} that turns predictive
uncertainty into prediction intervals with coverage guarantees and a principled
abstention threshold; (3) an evaluation protocol centred on calibration and
behaviour under distribution shift
(AffectNet~\cite{mollahosseini2017affectnet}$\rightarrow$AFEW-VA~\cite{kossaifi2017afew});
and (4) an on-device demonstrator that surfaces
the uncertainty and abstention to users without transmitting any imagery.

\paragraph{Scope and ethics.} The system estimates the user's own state at their
request. It is not a diagnostic tool and is not intended for assessing other
people. All inference in the demonstrator runs in the browser; no frames or
biometric signals leave the device.
